% !TEX program = xelatex
% ============================================================
%  chalk-annotate 示例文档 | Example Document
%  编译: xelatex example.tex (编译两次)
% ============================================================
\documentclass[aspectratio=169, 11pt]{beamer}

% ---- 中文支持 ----
\usepackage{ctex}

% ---- 手绘标注包 ----
\usepackage{chalk-annotate}

% ---- 常用包 ----
\usepackage{amsmath, amssymb}
\usepackage{booktabs}

% ============================================================
%  Beamer 主题
% ============================================================
\usetheme{default}
\setbeamertemplate{navigation symbols}{}
\setbeamertemplate{footline}{%
  \hfill\insertframenumber{}/\inserttotalframenumber\hspace{1em}\vspace{0.5em}%
}

\definecolor{accent}{HTML}{C0392B}
\setbeamercolor{frametitle}{fg=black}
\setbeamercolor{itemize item}{fg=accent}
\setbeamerfont{frametitle}{size=\large, series=\bfseries}

% ---- frametitle 样式 ----
\setbeamertemplate{frametitle}{%
  \vspace{0.5em}\insertframetitle\par
  \vspace{-0.2em}{\color{accent}\rule{2cm}{2pt}}\vspace{0.2em}%
}

% ============================================================
\begin{document}

% ---- 封面 ----
\begin{frame}[c]
  \centering
  {\Large\bfseries chalk-annotate}\par
  \vspace{0.5em}
  {\normalsize Hand-drawn chalk annotations for \LaTeX{} Beamer}\par
  \vspace{0.5em}
  {\small\color{gray} v1.0 \quad 2026-04-24}
\end{frame}

% ---- 基本用法 ----
\begin{frame}{Basic Usage}
  All commands take an optional \texttt{[color]} argument:

  \vspace{1em}
  \begin{itemize}\setlength{\itemsep}{0.6em}
    \item \texttt{\textbackslash hlellipse\{text\}}
          \quad$\Rightarrow$\quad \hlellipse{text}
    \item \texttt{\textbackslash hlunderline\{text\}}
          \quad$\Rightarrow$\quad \hlunderline{text}
    \item \texttt{\textbackslash hlbox\{text\}}
          \quad$\Rightarrow$\quad \hlbox{text}
    \item \texttt{\textbackslash hlarrowright\{text\}}
          \quad$\Rightarrow$\quad \hlarrowright{text}
    \item \texttt{\textbackslash hlstrike\{text\}}
          \quad$\Rightarrow$\quad \hlstrike{text}
  \end{itemize}
\end{frame}

% ---- 混合颜色 ----
\begin{frame}{Mixed Colors}
  Use different colors in the same document:

  \vspace{1em}
  \begin{itemize}\setlength{\itemsep}{0.8em}
    \item \hlellipse{Default (red)} — \texttt{\textbackslash hlellipse\{...\}}
    \item \hlellipse[blue]{Blue} — \texttt{\textbackslash hlellipse[blue]\{...\}}
    \item \hlellipse[green]{Green} — \texttt{\textbackslash hlellipse[green]\{...\}}
    \item \hlellipse[orange]{Orange} — \texttt{\textbackslash hlellipse[orange]\{...\}}
    \item \hlellipse[purple]{Purple} — \texttt{\textbackslash hlellipse[purple]\{...\}}
  \end{itemize}
\end{frame}

% ---- 中文示例 ----
\begin{frame}{中文示例}
  \begin{itemize}\setlength{\itemsep}{0.6em}
    \item 核心概念：\hlellipse{比较优势}
    \item 传统理论忽略了\hlunderline[blue]{要素禀赋结构}的内生性
    \item 实证研究发现 \hlbox[green]{$\beta_1 < 0$}：偏离比较优势抑制增长
    \item 对\hlellipse[orange]{中低收入国家}的影响更为显著
  \end{itemize}

  \vspace{1em}
  公式中使用：
  \[
    Growth_{it} = \beta_0
      + \hlellipse[blue]{$\beta_1$} \cdot TCI_{it}
      + \beta_2 \cdot X_{it}
      + \mu_i + \lambda_t + \varepsilon_{it}
  \]
\end{frame}

% ---- 方框展示 ----
\begin{frame}{Box Styles}
  \centering
  \vspace{1em}
  \renewcommand{\arraystretch}{2.5}
  \begin{tabular}{cl}
    \toprule
    \textbf{Color} & \textbf{Box with fill + border} \\
    \midrule
    \texttt{red}
      & \hlbox{示例文字 Sample text} \\
    \texttt{blue}
      & \hlbox[blue]{示例文字 Sample text} \\
    \texttt{green}
      & \hlbox[green]{示例文字 Sample text} \\
    \texttt{orange}
      & \hlbox[orange]{示例文字 Sample text} \\
    \texttt{purple}
      & \hlbox[purple]{示例文字 Sample text} \\
    \bottomrule
  \end{tabular}
\end{frame}

\end{document}
